% ======================================================================
% CHAPTER 3: MATHEMATICAL MODELING, NETWORK FORMALISM & PROBLEM FORMULATION
% ======================================================================
\chapter{Mathematical Modeling \& Problem Formulation}
\label{chap:formulation}

\section{Network Topology and Multigraph Mathematical Formalism}
We model a software-defined communication network as a directed multigraph $G = (V, E)$, where $V = \{1, 2, \dots, |V|\}$ represents the finite set of OpenFlow switching elements, and $E \subseteq V \times V$ represents the set of directed physical transmission links. Multiple directed edges may exist between adjacent switching pairs to model redundant trunk aggregates or bundled optical channels.

For each directed edge $e = (u, v) \in E$, the following physical and operational invariants are defined:
\begin{enumerate}
    \item \textbf{Transmission Bandwidth Capacity ($C_e \in \mathbb{R}^+$)}: The maximum line rate of the link in Megabits per second (Mbps). For bidirectional communication, links $e = (u, v)$ and $\bar{e} = (v, u)$ possess symmetric capacity $C_e = C_{\bar{e}}$, though their real-time traffic utilizations are independent.
    \item \textbf{Base Physical Delay ($d_{\text{base}}(e) \in \mathbb{R}^+$)}: The minimum, non-congested transit delay across link $e$, combining physical medium propagation delay and serializer transmission delay for standard Maximum Transmission Unit (MTU = 1500 bytes) frames:
    \begin{equation}
        d_{\text{base}}(e) = d_{\text{prop}}(e) + d_{\text{tx}}(e) = \frac{L_e}{v_{\text{med}}} + \frac{1500 \times 8}{C_e \times 10^6} \quad [\text{seconds}]
    \end{equation}
    where $L_e$ is the fiber/cable physical length in meters and $v_{\text{med}} \approx 2 \times 10^8\text{ m/s}$.
    \item \textbf{Instantaneous Link Utilization ($u_e(t) \in [0.0, 1.0]$)}: The dimensionless fraction of link capacity utilized by aggregate transmitted flows at time $t$:
    \begin{equation}
        u_e(t) = \min\left(1.0, \, \frac{\lambda_e(t)}{C_e}\right)
    \end{equation}
    where $\lambda_e(t)$ is the aggregate traffic arrival rate on link $e$ in Mbps.
\end{enumerate}

\section{The Multi-Commodity Flow Optimization Problem}

\subsection{Continuous Linear Formulation}
Let $\mathcal{K} = \{1, 2, \dots, K\}$ denote the active set of traffic commodities traversing network $G$. Each commodity $k \in \mathcal{K}$ is defined by a source ingress switch $s_k \in V$, a destination egress switch $t_k \in V$, and an instantaneous traffic demand $D_k > 0$.

If traffic flows are infinitely divisible and multi-path splitting is unrestricted, the optimal routing policy solves the classical fractional Multi-Commodity Flow (MCF) optimization problem:
\begin{align}
    \text{Minimize} \quad & \Psi(f) = \sum_{e \in E} \Phi_e\left( \frac{\sum_{k \in \mathcal{K}} f_e^k}{C_e} \right) \label{eq:mcf_objective} \\
    \text{subject to} \quad & \sum_{e \in \delta^+(v)} f_e^k - \sum_{e \in \delta^-(v)} f_e^k = b_v^k \quad \forall v \in V, \, \forall k \in \mathcal{K} \label{eq:mcf_conservation} \\
    & \sum_{k \in \mathcal{K}} f_e^k \le C_e \quad \forall e \in E \label{eq:mcf_capacity} \\
    & f_e^k \ge 0 \quad \forall e \in E, \, \forall k \in \mathcal{K} \label{eq:mcf_nonnegative}
\end{align}
where $f_e^k$ represents the flow rate of commodity $k$ traversing directed edge $e$, $\delta^+(v)$ and $\delta^-(v)$ denote the sets of outgoing and incoming edges incident on vertex $v$, and the divergence vector $b_v^k$ is:
\begin{equation}
    b_v^k = \begin{cases}
        D_k & \text{if } v = s_k \\
        -D_k & \text{if } v = t_k \\
        0 & \text{if } v \in V \setminus \{s_k, t_k\}
    \end{cases}
\end{equation}
The penalty function $\Phi_e(u)$ is a continuous, strictly convex function that models packet queuing delay and buffer overflow risk as link utilization approaches unity.

\subsection{Unsplittable Flow Integer Formulation and Complexity}
In packet-switched computer networks, splitting a single transport-layer TCP flow across multiple paths with disparate latencies induces severe packet reordering. Out-of-order packet arrivals trigger TCP Fast Retransmit and cause the TCP sender to artificially halve its congestion window ($cwnd$), collapsing application throughput.

Consequently, production traffic engineering mandates \textbf{Unsplittable Multi-Commodity Flow (UMCF)}, wherein each commodity $k$ must follow a single, atomic path from source to destination. We introduce binary decision variables:
\begin{equation}
    x_e^k \in \{0, 1\} \quad \forall e \in E, \, \forall k \in \mathcal{K}
\end{equation}
where $x_e^k = 1$ if commodity $k$ traverses edge $e$, and $0$ otherwise. The flow conservation constraint becomes:
\begin{equation}
    \sum_{e \in \delta^+(v)} x_e^k - \sum_{e \in \delta^-(v)} x_e^k = \begin{cases}
        1 & \text{if } v = s_k \\
        -1 & \text{if } v = t_k \\
        0 & \text{otherwise}
    \end{cases} \quad \forall v \in V, \, \forall k \in \mathcal{K}
\end{equation}
and the link capacity constraint becomes:
\begin{equation}
    \sum_{k \in \mathcal{K}} D_k \cdot x_e^k \le C_e \quad \forall e \in E
\end{equation}

\begin{theorem}[$\mathcal{NP}$-Hardness of Unsplittable Traffic Engineering]
The Unsplittable Multi-Commodity Flow problem with capacity constraints is strictly $\mathcal{NP}$-hard.
\end{theorem}

\begin{proof}
We prove $\mathcal{NP}$-hardness by polynomial-time reduction from the classical \textsc{2-Directed Disjoint Paths (2-DDP)} problem, which is proven to be $\mathcal{NP}$-complete \cite{garey1979computers}.

An instance of 2-DDP consists of a directed graph $G = (V, E)$ and two source-terminal pairs $(s_1, t_1)$ and $(s_2, t_2)$. The question is whether there exist two edge-disjoint paths $\mathcal{P}_1$ (from $s_1$ to $t_1$) and $\mathcal{P}_2$ (from $s_2$ to $t_2$) such that $E(\mathcal{P}_1) \cap E(\mathcal{P}_2) = \emptyset$.

We transform this 2-DDP instance into an instance of UMCF:
\begin{enumerate}
    \item Assign unit capacity to all edges: $C_e = 1$ for all $e \in E$.
    \item Define two commodities $\mathcal{K} = \{1, 2\}$ with demands $D_1 = 1$ and $D_2 = 1$.
    \item Set the source-destination pairs as $(s_1, t_1)$ and $(s_2, t_2)$.
\end{enumerate}
In any valid solution to this UMCF instance, for each edge $e \in E$, the capacity constraint requires:
\begin{equation}
    D_1 \cdot x_e^1 + D_2 \cdot x_e^2 = 1 \cdot x_e^1 + 1 \cdot x_e^2 \le 1
\end{equation}
Because $x_e^k \in \{0, 1\}$, it is impossible for both $x_e^1 = 1$ and $x_e^2 = 1$ on any edge $e$. Therefore, the paths selected for commodity 1 and commodity 2 must be strictly edge-disjoint. Thus, a feasible unsplittable flow exists if and only if the 2-DDP instance has a solution. Since the reduction takes linear time, UMCF is $\mathcal{NP}$-hard.
\end{proof}

Because UMCF is $\mathcal{NP}$-hard, computing exact global integer linear programs in real-time under sub-second traffic variations is computationally intractable. This theoretical bottleneck establishes the imperative for an autonomous Deep Reinforcement Learning policy capable of sub-millisecond heuristic inference.

\section{Markov Decision Process (MDP) Formulation}
When a new flow $f = (s, t, \text{demand})$ arrives at ingress switch $s$ destined for egress switch $t$, the SDN controller intercepts the first packet miss and invokes the DRL routing agent. We formalize this decision loop as an MDP:

\subsection{State Space Representation ($\mathcal{S} \in [0.0, 1.0]^{10}$)}
A major barrier in applying reinforcement learning to networking is achieving topological invariance. If the state vector directly encodes raw node IDs or per-link metrics, the neural network structure becomes permanently tied to a specific switch count $|V|$ and link count $|E|$. 

To solve this, we introduce a \textbf{Topologically Normalized 10-Dimensional State Vector} $s_t \in \mathbb{R}^{10}$:
\begin{equation}
    s_t = \begin{bmatrix} 
        s_0 \\ s_1 \\ s_2 \\ s_3 \\ s_4 \\ s_5 \\ s_6 \\ s_7 \\ s_8 \\ s_9 
    \end{bmatrix} = \begin{bmatrix}
        \min\left(1.0, \, \max\left(0.1, \, \dfrac{|\mathcal{P}_0| - 1}{10.0}\right)\right) \\[0.5em]
        \max_{e \in \mathcal{P}_0} u(e) \\[0.5em]
        \max_{e \in \mathcal{P}_1} u(e) \\[0.5em]
        \max_{e \in \mathcal{P}_2} u(e) \\[0.5em]
        \max_{e \in \mathcal{P}_3} u(e) \\[0.5em]
        \min\left(1.0, \, \dfrac{\sum_{e \in \mathcal{P}_0} d_{\text{base}}(e)}{50.0}\right) \\[0.5em]
        \min\left(1.0, \, \dfrac{\sum_{e \in \mathcal{P}_1} d_{\text{base}}(e)}{50.0}\right) \\[0.5em]
        \min\left(1.0, \, \dfrac{\sum_{e \in \mathcal{P}_2} d_{\text{base}}(e)}{50.0}\right) \\[0.5em]
        \min\left(1.0, \, \dfrac{\sum_{e \in \mathcal{P}_3} d_{\text{base}}(e)}{50.0}\right) \\[0.5em]
        \max_{e \in E} u(e)
    \end{bmatrix}
    \label{eq:state_vector_full}
\end{equation}

Table~\ref{tab:state_features_breakdown} details the physical interpretation, mathematical domain, and architectural rationale for each feature.

\begin{table}[htbp]
\centering
\caption{Detailed Mathematical Specification of the 10-Dimensional State Space Vector}
\label{tab:state_features_breakdown}
\small
\begin{tabularx}{\textwidth}{c l c X}
\toprule
\textbf{Feature} & \textbf{Mathematical Definition} & \textbf{Domain} & \textbf{Physical Interpretation \& Role} \\
\midrule
$s_0$ & $\text{clamp}\left(\frac{|\mathcal{P}_0| - 1}{10.0}, 0.1, 1.0\right)$ & $[0.1, 1.0]$ & Normalized hop distance of primary shortest path. Encodes spatial topological diameter. \\
\addlinespace
$s_1$ & $\max_{e \in \mathcal{P}_0} u(e)$ & $[0.0, 1.0]$ & Peak bottleneck link utilization along Candidate Path 0 (Shortest Path). \\
\addlinespace
$s_2$ & $\max_{e \in \mathcal{P}_1} u(e)$ & $[0.0, 1.0]$ & Peak bottleneck link utilization along Candidate Path 1 (Second Shortest Path). \\
\addlinespace
$s_3$ & $\max_{e \in \mathcal{P}_2} u(e)$ & $[0.0, 1.0]$ & Peak bottleneck link utilization along Candidate Path 2 (Edge-Disjoint Bypass). \\
\addlinespace
$s_4$ & $\max_{e \in \mathcal{P}_3} u(e)$ & $[0.0, 1.0]$ & Peak bottleneck link utilization along Candidate Path 3 (Node-Diverse / Widest Path). \\
\addlinespace
$s_5$ & $\min\left(1.0, \frac{\sum_{e \in \mathcal{P}_0} d_{\text{base}}(e)}{50.0}\right)$ & $[0.0, 1.0]$ & Cumulative propagation latency of Path 0, normalized against 50~ms scale. \\
\addlinespace
$s_6$ & $\min\left(1.0, \frac{\sum_{e \in \mathcal{P}_1} d_{\text{base}}(e)}{50.0}\right)$ & $[0.0, 1.0]$ & Cumulative propagation latency of Path 1, normalized against 50~ms scale. \\
\addlinespace
$s_7$ & $\min\left(1.0, \frac{\sum_{e \in \mathcal{P}_2} d_{\text{base}}(e)}{50.0}\right)$ & $[0.0, 1.0]$ & Cumulative propagation latency of Path 2, normalized against 50~ms scale. \\
\addlinespace
$s_8$ & $\min\left(1.0, \frac{\sum_{e \in \mathcal{P}_3} d_{\text{base}}(e)}{50.0}\right)$ & $[0.0, 1.0]$ & Cumulative propagation latency of Path 3, normalized against 50~ms scale. \\
\addlinespace
$s_9$ & $\max_{e \in E} u(e)$ & $[0.0, 1.0]$ & Network-wide global peak link utilization. Encodes overall network stress. \\
\bottomrule
\end{tabularx}
\end{table}

\begin{theorem}[Topological Invariance of State Vector $s_t$]
The state representation $s_t$ is invariant to graph isomorphic relabeling, switch permutations, and absolute graph cardinality $|V|$ and $|E|$.
\end{theorem}

\begin{proof}
Let $G_1 = (V_1, E_1)$ and $G_2 = (V_2, E_2)$ be two arbitrary network graphs. Feature $s_0$ depends strictly on scalar path length $|\mathcal{P}_0|$, scaled by constant $10.0$. Features $s_1$--$s_4$ and $s_9$ depend strictly on scalar link utilization ratios $u(e) \in [0.0, 1.0]$, which are dimensionless quotients of throughput over capacity. Features $s_5$--$s_8$ depend on scalar physical latencies scaled by $50.0\text{ ms}$. None of the features index switch identifiers, port numbers, or matrix dimensions. Therefore, the vector dimensions and semantic meaning remain strictly identical regardless of graph scale, enabling zero-shot model evaluation across arbitrary graphs.
\end{proof}

\subsection{Action Space Formulation ($\mathcal{A} \in \{0, 1, 2, 3\}$)}
Rather than permitting the neural network to output unconstrained next-hop ports at each switch (which creates an exponentially scaling action space and high loop risk), our architecture uses an action space of size $K=4$ representing diversity-aware end-to-end candidate paths:
\begin{equation}
    \mathcal{A} = \{0, 1, 2, 3\}
\end{equation}

Candidate paths are synthesized via Yen's $K$-Shortest Paths algorithm \cite{yen1971finding} augmented with diversity filtering:
\begin{enumerate}
    \item \textbf{Action 0 ($\mathcal{P}_0$)}: Standard Dijkstra Shortest Path (minimum hop count and minimal propagation delay).
    \item \textbf{Action 1 ($\mathcal{P}_1$)}: Second Shortest Simple Path.
    \item \textbf{Action 2 ($\mathcal{P}_2$)}: Edge-Diverse Lateral Bypass Path. Synthesized by selecting the simple path that minimizes edge overlap with $\mathcal{P}_0$:
    \begin{equation}
        \mathcal{P}_2 = \arg\min_{\mathcal{P} \in \mathcal{P}_{\text{cand}} \setminus \{\mathcal{P}_0, \mathcal{P}_1\}} \sum_{e \in \mathcal{P}} \mathbb{I}(e \in \mathcal{P}_0)
    \end{equation}
    \item \textbf{Action 3 ($\mathcal{P}_3$)}: Node-Diverse / Widest Path. Synthesized by selecting the path that maximizes residual bottleneck bandwidth across all intermediate transit nodes disjoint from $\mathcal{P}_0$:
    \begin{equation}
        \mathcal{P}_3 = \arg\max_{\mathcal{P} \in \mathcal{P}_{\text{cand}} \setminus \{\mathcal{P}_0, \mathcal{P}_1, \mathcal{P}_2\}} \min_{e \in \mathcal{P}} \left( C(e) - \text{Throughput}(e) \right)
    \end{equation}
\end{enumerate}

\subsection{Multi-Objective Reward Function}
The scalar reward function balances five competing objectives: hop overhead, propagation and queuing latency, network jitter, packet loss rate, and peak link congestion:
\begin{equation}
    \mathcal{R}(s_t, a_t) = - \Big( w_{\text{hop}} \cdot |\mathcal{P}_a| + w_{\text{lat}} \cdot D(\mathcal{P}_a) + \Phi_{\text{cong}}(U_{\max}) + w_{\text{jit}} \cdot J(\mathcal{P}_a) + w_{\text{loss}} \cdot P_{\text{loss}}(\mathcal{P}_a) \Big)
    \label{eq:multi_objective_reward}
\end{equation}
where $U_{\max} = \max_{e \in \mathcal{P}_a} u(e)$ is the peak bottleneck link utilization on the chosen path, and the calibrated weights are:
\begin{equation}
    w_{\text{hop}} = 0.35, \quad w_{\text{lat}} = 0.06, \quad w_{\text{jit}} = 0.25, \quad w_{\text{loss}} = 0.50
\end{equation}

\section{Derivation of the Asymptotic Congestion Barrier Penalty}
A pivotal theoretical contribution of this project is the formulation of the \textbf{Asymptotic Congestion Barrier Penalty} $\Phi_{\text{cong}}(U_{\max})$. In classical interior-point optimization, barrier functions (such as the Frisch logarithmic barrier $-\mu \ln(C - f)$) prevent optimization variables from violating feasibility boundaries by generating infinite penalty barriers at constraints.

We construct a specialized piecewise continuous barrier function:
\begin{equation}
    \Phi_{\text{cong}}(U_{\max}) = \begin{cases}
        1.5 \cdot U_{\max} & \text{if } U_{\max} \le 0.70 \\[0.5em]
        12.0 \cdot \left( \dfrac{U_{\max}^{1.8}}{\max(0.01, \, 1.02 - U_{\max})} \right) & \text{if } U_{\max} > 0.70
    \end{cases}
    \label{eq:barrier_penalty_formal}
\end{equation}

\subsection{Mathematical Analysis and Gradient Dynamics}
We analyze the first-order derivative of $\Phi_{\text{cong}}$ with respect to bottleneck utilization $u = U_{\max}$ in the barrier regime ($u > 0.70$):
\begin{equation}
    \Phi_{\text{cong}}(u) = 12.0 \cdot \frac{u^{1.8}}{1.02 - u}
\end{equation}
Applying the quotient rule:
\begin{align}
    \frac{d \Phi_{\text{cong}}}{du} &= 12.0 \cdot \frac{1.8 u^{0.8}(1.02 - u) - u^{1.8}(-1)}{(1.02 - u)^2} \nonumber \\
    &= 12.0 \cdot \frac{1.836 u^{0.8} - 1.8 u^{1.8} + u^{1.8}}{(1.02 - u)^2} \nonumber \\
    &= 12.0 \cdot \frac{1.836 u^{0.8} - 0.8 u^{1.8}}{(1.02 - u)^2} > 0 \quad \forall u \in (0.70, 1.0]
\end{align}
Furthermore, the second derivative satisfies:
\begin{equation}
    \frac{d^2 \Phi_{\text{cong}}}{du^2} > 0 \quad \forall u \in (0.70, 1.0]
\end{equation}
confirming that $\Phi_{\text{cong}}(u)$ is strictly convex and monotonically escalating.

\begin{table}[htbp]
\centering
\caption{Numerical Escalation of the Asymptotic Barrier Penalty vs. Linear Hop Penalty}
\label{tab:barrier_numerical_values}
\begin{tabular}{c c c c l}
\toprule
\textbf{Utilization $U_{\max}$} & \textbf{Regime} & \textbf{Barrier Penalty $\Phi_{\text{cong}}$} & \textbf{Hop Penalty ($0.35 \cdot 3$ hops)} & \textbf{Dominant Agent Incentive} \\
\midrule
0.20 & Linear & 0.30 & 1.05 & Conserve hops; select shortest path $\mathcal{P}_0$ \\
0.40 & Linear & 0.60 & 1.05 & Conserve hops; select shortest path $\mathcal{P}_0$ \\
0.60 & Linear & 0.90 & 1.05 & Conserve hops; select shortest path $\mathcal{P}_0$ \\
0.70 & Boundary & 1.05 & 1.05 & Indifference threshold between hop count and load \\
0.75 & Barrier & 26.23 & 1.05 & \textbf{Strong incentive to bypass core} \\
0.85 & Barrier & 53.31 & 1.05 & \textbf{Aggressive detour enforcement} \\
0.95 & Barrier & 156.42 & 1.05 & \textbf{Absolute emergency rerouting} \\
0.99 & Barrier & 386.11 & 1.05 & \textbf{Asymptotic penalty saturation} \\
\bottomrule
\end{tabular}
\end{table}

As proven in Table~\ref{tab:barrier_numerical_values}, below $70\%$ load, the hop penalty dominates, forcing the agent to pick shortest paths and conserve global network resources. However, once a core link exceeds $70\%$ utilization, the barrier penalty escalates to $53.31$ and $156.42$, utterly dwarfing the modest 1.05 hop penalty. This non-linear mathematical gradient forces the agent to steer traffic onto lateral bypass paths before physical switch buffer dropouts can manifest.

\section{Analytical Queuing, Jitter, and Packet Loss Modeling}

\subsection{Path Latency Formulation via Kleinrock's Approximation}
The end-to-end latency $D(\mathcal{P})$ across path $\mathcal{P}$ consists of three additive delay components on each traversed link $e$:
\begin{equation}
    D(\mathcal{P}) = \sum_{e \in \mathcal{P}} \Big( d_{\text{prop}}(e) + d_{\text{tx}}(e) + d_{\text{queue}}(e, u_e) \Big)
\end{equation}
Under Kleinrock's network independence approximation \cite{kleinrock1964communication}, packet arrivals at each internal switch output queue follow Poisson distributions with average rate $\lambda_e$, and packet service times are exponentially distributed with mean service rate $\mu_e = C_e / \bar{L}$, where $\bar{L}$ is the average packet size.

Each link buffer operates as an $M/M/1$ queue. The steady-state average queuing delay is given by:
\begin{equation}
    W_q(e) = \frac{\lambda_e}{\mu_e(\mu_e - \lambda_e)} = \frac{u_e}{\mu_e(1 - u_e)}
\end{equation}
To avoid numerical singularity as $u_e \to 1.0$ and model practical buffer limiters, we formulate the calibrated queuing delay model:
\begin{equation}
    d_{\text{queue}}(e, u_e) = d_{\text{base}}(e) \cdot \left( 0.2 + \frac{0.8}{\max\left(0.05, \, 1.0 - \min(0.95, \, u_e)\right)} \right)
    \label{eq:queuing_delay_formula}
\end{equation}
Under zero load ($u_e = 0$), $d_{\text{queue}} = d_{\text{base}}$. When load reaches $95\%$, the denominator collapses to $0.05$, scaling queuing delay by sixteen-fold to $16.2 \cdot d_{\text{base}}$, accurately modeling real-world buffer congestion.

\subsection{Network Jitter Model (RFC 3393 / RFC 3550)}
Packet delay variation (jitter) measures the statistical variance in one-way end-to-end transit times. Standardized in RFC 3393 \cite{deroo2002ip} and RFC 3550 \cite{schulzrinne2003rtp}, jitter is evaluated via an Exponential Moving Average (EMA) filter:
\begin{equation}
    J_t = J_{t-1} + \frac{|D_t - D_{t-1}| - J_{t-1}}{16}
\end{equation}
In our analytical network queuing engine, jitter volatility is modeled directly as a function of peak path utilization $U_{\max}$:
\begin{equation}
    J(\mathcal{P}) = 0.35 \cdot |\mathcal{P}| + 0.45 \cdot D(\mathcal{P}) \cdot \min\left(5.0, \, \frac{U_{\max}^2}{\max(0.02, \, 1.0 - \min(0.98, \, U_{\max}))}\right)
    \label{eq:jitter_analytical_model}
\end{equation}
When links operate under low load, queue sizes are constant, and jitter remains minimal ($< 1\text{ ms}$). As utilization surpasses $80\%$, bursty packet arrivals cause queue occupancy to oscillate violently, inducing massive jitter spikes.

\subsection{Buffer Overflow Packet Loss Model ($M/M/1/K$)}
In physical OpenFlow switches, output port queues possess a finite buffer capacity of $K$ packets (typically 100 to 1000 packets in commodity hardware). Under the $M/M/1/K$ queuing model, packet loss occurs when an incoming packet arrives while all $K$ buffer slots are occupied.

The stationary distribution of an $M/M/1/K$ queue with traffic intensity $\rho = \lambda / \mu$ yields the blocking probability $P_K$:
\begin{equation}
    P_{\text{loss}} = P_K = \frac{(1 - \rho)\rho^K}{1 - \rho^{K+1}}
\end{equation}
When traffic intensity $\rho = u \le 0.70$ and buffer size $K \ge 100$, the theoretical loss probability is astronomically small ($P_{\text{loss}} < 10^{-6}$). However, when utilization surges beyond $80\%$, microbursts exhaust the remaining buffer slots, triggering drop-tail packet loss. We formulate this empirical reality via:
\begin{equation}
    P_{\text{loss}}(u) = \begin{cases}
        0.01\% & \text{if } u \le 0.70 \\[0.5em]
        0.01\% + 35.0 \cdot \left( \dfrac{u - 0.70}{0.30} \right)^{2.2} & \text{if } u > 0.70
    \end{cases}
    \label{eq:packet_loss_model}
\end{equation}
At $70\%$ load, packet loss is $0.01\%$. At $85\%$ load, loss rises to $7.6\%$. At $100\%$ saturation, packet loss escalates to $35.01\%$, reflecting severe buffer overflow.

\subsection{Jain's Fairness Index Formulation}
To rigorously quantify the global balance of link loads across the entire network topology, we compute Jain's Fairness Index \cite{jain1984quantitative}:
\begin{equation}
    \mathcal{J}(u) = \frac{\left( \sum_{e \in E} u(e) \right)^2}{|E| \cdot \sum_{e \in E} u(e)^2}
    \label{eq:jains_fairness_formula}
\end{equation}

\begin{theorem}[Properties of Jain's Fairness Index]
Jain's Fairness Index satisfies:
\begin{enumerate}
    \item \textbf{Boundedness}: $\frac{1}{|E|} \le \mathcal{J}(u) \le 1.0$.
    \item \textbf{Scale Invariance}: $\mathcal{J}(\alpha u) = \mathcal{J}(u)$ for all scalar multipliers $\alpha > 0$.
    \item \textbf{Continuity}: $\mathcal{J}(u)$ is smooth and continuous for all non-zero link load vectors.
\end{enumerate}
\end{theorem}

\begin{proof}
Let $u = [u_1, \dots, u_{|E|}]^T \in \mathbb{R}^{|E|}$ and let $\mathbf{1} = [1, \dots, 1]^T$. By the Cauchy--Schwarz inequality:
\begin{equation}
    \left( \sum_{e \in E} u(e) \cdot 1 \right)^2 \le \left( \sum_{e \in E} u(e)^2 \right) \left( \sum_{e \in E} 1^2 \right) = |E| \sum_{e \in E} u(e)^2
\end{equation}
Dividing both sides by $|E| \sum_{e \in E} u(e)^2$ yields $\mathcal{J}(u) \le 1.0$. Equality holds if and only if $u(e)$ is identical across all links ($u_1 = u_2 = \dots = u_{|E|}$).

Conversely, if all traffic is concentrated onto a single link $e_1$ such that $u(e_1) = U$ and $u(e_i) = 0$ for all $i \ge 2$, then:
\begin{equation}
    \mathcal{J}(u) = \frac{(U)^2}{|E| \cdot (U^2)} = \frac{1}{|E|}
\end{equation}
Scale invariance follows immediately from homogeneity of degree 2 in both numerator and denominator.
\end{proof}

A high fairness index ($\mathcal{J} \to 1.0$) demonstrates that traffic is balanced evenly across the entire mesh fabric, avoiding destructive core hotspots.
